\documentclass[11pt]{article}
\usepackage[margin=1in]{geometry}
\usepackage{amsmath,amssymb,amsthm}
\usepackage{xurl}
\usepackage{hyperref}
\sloppy
\let\oldus\_ \renewcommand{\_}{\oldus\allowbreak}
\newtheorem{theorem}{Theorem}
\newtheorem{lemma}[theorem]{Lemma}
\theoremstyle{definition}
\newtheorem{definition}[theorem]{Definition}

\title{Disproof of conjecture 00000003238:\\ a direct sum of nilpotent Jordan blocks has the closed
unit disc as spectrum, while every block has spectrum $\{0\}$}
\author{Jackmeson1}
\date{2026-10-06}

\begin{document}
\maketitle

\section{The conjecture and the clause refuted}

\textbf{Statement (English, verbatim).} Definition: Edge spectra: the block spectral structure of
direct-sum operators. Conjecture: The edge theorem: the closure of the spectrum of a direct sum of
operators is the union closure of block spectra, with limit points the accumulated edges of the
family of summands. (direct sum block spectrum accumulated edge) (The Chinese text states the
same.)

\textbf{Reading.} Let $(H_n)_{n\ge 0}$ be complex Hilbert spaces and $(A_n)$ a uniformly bounded
family of bounded operators $A_n:H_n\to H_n$. Their direct sum is the operator
$\bigoplus_n A_n$ on the Hilbert direct sum $\bigoplus_n H_n$ (square-summable sequences
$x=(x_n)$, $x_n\in H_n$) given by $(\bigoplus_n A_n)x=(A_nx_n)_n$. The conjecture's main clause is
\[
  \text{(E)}\qquad \overline{\sigma\Big(\bigoplus_n A_n\Big)}
     \;=\;\overline{\bigcup_n \sigma(A_n)} .
\]
The clause ``with limit points the accumulated edges of the family of summands'' is joined to (E)
by ``with'' (Chinese: ``and''); it uses the undefined term ``accumulated edges'' and is not needed:
refuting (E) refutes the conjecture. We refute (E) for one explicit family of contractions
($\|A_n\|\le 1$) on finite-dimensional spaces $H_n=\mathbb{C}^{n+1}$.

\textbf{Readings not refuted.} The phrases ``union closure'', ``limit points'' and ``family of
summands'' indicate an infinite family, which is what we refute. For a \emph{finite} direct sum
the identity (E) is true, and it is also true for families of \emph{normal} operators (for normal
$A_n$ and $\lambda$ at positive distance $d$ from $\bigcup_n\sigma(A_n)$, the resolvents satisfy
$\|(A_n-\lambda)^{-1}\|\le 1/d$, so they assemble into a bounded inverse). Neither restriction
appears in the text, which says ``a direct sum of operators'' with no normality hypothesis. We
also only treat the Hilbert ($\ell^2$) direct sum, which is the standard direct sum of Hilbert
space operators.

\section{Definitions (as formalized)}

All conventions below are those of the Lean file.
\begin{itemize}
\item $\mathbb{N}=\{0,1,2,\dots\}$. Block $n$ is $\mathbb{C}^{n+1}$ with the Euclidean norm
  (\texttt{Block n := EuclideanSpace C (Fin (n+1))}), coordinates indexed by $0,\dots,n$.
\item The Hilbert direct sum of a family $E:\mathbb{N}\to\text{Type}$ of complex inner product
  spaces is Mathlib's \texttt{lp E 2}, the square-summable sections, with Mathlib's inner product
  structure \texttt{lp.instInnerProductSpace} (Mathlib calls it the ``Hilbert sum'' of the
  family). For our blocks it is complete (a Hilbert space).
\item \texttt{dsum A hA}: for $A_n:E_n\to E_n$ bounded with $\|A_n\|\le 1$, the bounded operator
  on \texttt{lp E 2} with $(\texttt{dsum A hA}\,x)_n = A_n x_n$ (theorem \texttt{dsum\_apply}); it
  maps the $n$-th summand by $A_n$ (\texttt{dsum\_single}), is the unique bounded operator acting
  blockwise as $(A_n)$ (\texttt{dsum\_unique}), and $\|\texttt{dsum A hA}\|\le 1$
  (\texttt{norm\_dsum\_le}).
\item Spectrum: Mathlib's \texttt{spectrum C a} $=\{\mu\in\mathbb{C} : \mu\cdot 1-a \text{ is not a
  unit}\}$, computed in the Banach algebra of bounded (continuous) linear operators. Being a unit
  there means having a bounded two-sided inverse.
\item The shift $J_n$ on $\mathbb{C}^{n+1}$: $(J_nv)_j=v_{j+1}$ for $j<n$ and $(J_nv)_n=0$
  (\texttt{shift n}). Equivalently $J_ne_0=0$, $J_ne_{j+1}=e_j$: a single nilpotent Jordan block of
  size $n+1$.
\item \texttt{EdgeTheorem}: the statement that (E) holds for every family $E:\mathbb{N}\to\text{Type}$
  of complex Hilbert spaces and every family $(A_n)$ with $\|A_n\|\le 1$. Contractive families are
  uniformly bounded, so \texttt{EdgeTheorem} is a special case of the conjecture's claim.
\end{itemize}

\section{Proof}

\begin{lemma}[blocks]\label{lem:blocks}
$\|J_n\|\le 1$, $J_n^{\,n+1}=0$, and $\sigma(J_n)=\{0\}$. Hence
$\overline{\bigcup_n\sigma(J_n)}=\{0\}$.
\end{lemma}
\begin{proof}
$\|J_nv\|^2=\sum_{j<n}|v_{j+1}|^2\le\|v\|^2$. By induction on $k$,
$(J_n^kv)_j=v_{j+k}$ if $j+k\le n$ and $0$ otherwise, so $J_n^{n+1}=0$. If $\mu\ne0$, then
$\mu\cdot1-J_n$ is the sum of the unit $\mu\cdot 1$ and the nilpotent $-J_n$, which commute, so it is
a unit. For $\mu=0$, $-J_n$ is nilpotent in a nontrivial ring, so it is not a unit.
\end{proof}

Let $T=\bigoplus_n J_n$ (\texttt{T := dsum shift norm\_shift\_le}) on
$\bigoplus_n\mathbb{C}^{n+1}$; $\|T\|\le1$.

\begin{lemma}[test vectors]\label{lem:test}
For $z\in\mathbb{C}$ let $v^{(n)}=(1,z,z^2,\dots,z^n)\in\mathbb{C}^{n+1}$. Then
$(z\cdot1-J_n)v^{(n)}=z^{n+1}e_n$, and in the direct sum
$(z\cdot 1-T)\,\iota_n(v^{(n)})=\iota_n(z^{n+1}e_n)$, where $\iota_n$ is the inclusion of the
$n$-th summand.
\end{lemma}
\begin{proof}
For $j<n$ the $j$-th coordinate is $z\cdot z^j-z^{j+1}=0$; the $n$-th is $z\cdot z^n-0$. The second
identity follows since $T\iota_n=\iota_nJ_n$.
\end{proof}

\begin{theorem}\label{thm:disc}
$\sigma(T)$ is the closed unit disc. In particular $\tfrac12\in\sigma(T)$.
\end{theorem}
\begin{proof}
($\subseteq$) $\sigma(T)\subseteq\{|z|\le\|T\|\,\|1\|\}$ and $\|T\|\le1$, $\|1\|\le1$.
($\supseteq$) Let $|z|<1$ and suppose $z\cdot1-T$ had a bounded inverse $S$. Put
$x_n=\iota_n(v^{(n)})$. Then $\|x_n\|=\|v^{(n)}\|\ge|v^{(n)}_0|=1$, and by
Lemma~\ref{lem:test}
\[
 1\le\|x_n\|=\|S(z\cdot1-T)x_n\|\le\|S\|\,\|z^{n+1}e_n\|=\|S\|\,|z|^{n+1}\qquad(n\ge0).
\]
Choose $n$ with $|z|^n<1/(\|S\|+1)$; then $\|S\|\,|z|^{n+1}\le\|S\|\,|z|^n<1$, a contradiction.
So the open disc lies in $\sigma(T)$; since $\sigma(T)$ is closed, so does its closure, the closed
disc.
\end{proof}

\begin{theorem}[main]
$\overline{\sigma(\bigoplus_nJ_n)}\ne\overline{\bigcup_n\sigma(J_n)}$. Hence \texttt{EdgeTheorem}
is false, and the conjecture's clause (E) fails for the uniformly bounded family $(J_n)$.
\end{theorem}
\begin{proof}
By Lemma~\ref{lem:blocks} the right side is $\{0\}$; by Theorem~\ref{thm:disc} the left side
contains $\tfrac12\ne0$.
\end{proof}

\section{Lean formalization and axiom audit}

File \texttt{lean/Conjecture3238/Basic.lean} (273 lines, only \texttt{import Mathlib}, Lean and
Mathlib v4.33.1), namespace \texttt{C3238}. Main statements:
\begin{itemize}
\item \texttt{theorem not\_edgeTheorem : $\neg$ EdgeTheorem}
\item \texttt{theorem edge\_theorem\_counterexample : closure (spectrum C T) $\ne$ closure
  ($\bigcup$ n, spectrum C (shift n))}
\item \texttt{theorem spectrum\_T : spectrum C T = Metric.closedBall 0 1}
\item \texttt{theorem spectrum\_shift (n) : spectrum C (shift n) = \{0\}}
\item \texttt{theorem closure\_iUnion\_spectrum\_shift : closure ($\bigcup$ n, spectrum C (shift n))
  = \{0\}}
\item \texttt{theorem mem\_spectrum\_T \{z\} (hz : $\|$z$\|$ < 1) : z $\in$ spectrum C T}
\item supporting: \texttt{dsum\_apply}, \texttt{dsum\_single}, \texttt{dsum\_unique},
  \texttt{norm\_dsum\_le}, \texttt{norm\_shift\_le}, \texttt{shift\_pow\_eq\_zero},
  \texttt{block\_identity}, \texttt{sum\_identity}.
\end{itemize}
Here \texttt{C} stands for $\mathbb{C}$ and \texttt{T := dsum shift norm\_shift\_le}.

\textbf{Axiom audit.} \texttt{\#print axioms} for \texttt{not\_edgeTheorem},
\texttt{edge\_theorem\_counterexample}, \texttt{spectrum\_T}, \texttt{spectrum\_shift} and
\texttt{dsum\_unique} reports exactly \texttt{[propext, Classical.choice, Quot.sound]}. There is no
\texttt{sorry}, no \texttt{native\_decide}, and no added axiom.

\section*{Sources}
\begin{itemize}
\item Wikipedia, \emph{Spectrum (functional analysis)},
  \url{https://en.wikipedia.org/wiki/Spectrum_(functional_analysis)}: ``the spectrum of a bounded
  linear operator (or, more generally, an unbounded linear operator) is a generalisation of the set
  of eigenvalues of a matrix'' and ``More generally, by the bounded inverse theorem, T is not
  invertible if it is not bounded below''. (Background only; the Lean proof of
  $\tfrac12\in\sigma(T)$ is self-contained.)
\item Mathlib, \texttt{Mathlib/Analysis/InnerProductSpace/l2Space.lean} (definition of the Hilbert
  sum \texttt{lp G 2}) and \texttt{Mathlib/Algebra/Algebra/Spectrum/Basic.lean} (\texttt{spectrum}).
\end{itemize}

\end{document}
